% TOP_PROBLEM: {"id": 3060, "title": "The additive basis conjecture", "category_id": 2, "category": {"id": 2, "name": "combinatorics", "display_name": "Combinatorics", "description": "Counting problems, graph theory, discrete structures.", "slug": "combinatorics", "order_index": 2, "created_at": "2026-07-31T15:26:25.671Z"}, "status": "open", "problem_number": "OPG-150", "set_id": 12}
% FIRST_POSTED: 2026-10-01
% ENTRY_KIND: statement_only
% UnsolvedMath ID 3060; source catalogue code OPG-150
% !TeX program = lualatex
% Definition and sources only; this file is not an LLM solution attempt.
\documentclass[11pt,a4paper]{article}
\usepackage[margin=25mm]{geometry}
\usepackage{fontspec}
\setmainfont{lmroman10-regular.otf}[BoldFont=lmroman10-bold.otf,
 ItalicFont=lmroman10-italic.otf,BoldItalicFont=lmroman10-bolditalic.otf]
\usepackage{amsmath,amssymb,mathtools}
\usepackage{unicode-math}
\setmathfont{latinmodern-math.otf}
\AtBeginDocument{\let\setminus\smallsetminus}
\usepackage{xurl,hyperref}
\hypersetup{colorlinks=true,urlcolor=blue}
\setcounter{secnumdepth}{0}
\setlength{\parindent}{0pt}
\setlength{\parskip}{5pt}
\setlength{\emergencystretch}{3em}
\begin{document}
\section*{The additive basis conjecture}\label{3060}
{\small UnsolvedMath ID 3060 · OPG-150 · Combinatorics · OpenGarden}
\subsection{Definitions and mathematical statement}
For a prime $p$ and an integer $n\ge1$, let $V=\mathbb F_p^n$.
An additive basis here is a finite multiset $M$ of vectors for which every
$v\in V$ is the sum of a submultiset of $M$. Each occurrence can be used
at most once; equal vectors occurring several times remain separately
available. The empty sum is allowed and equals the zero vector.
This definition differs from an ordinary vector-space basis, where
arbitrary field coefficients are permitted.

The conjecture asks for a positive integer $c(p)$ depending only on $p$
such that, for every $n$ and every sequence of vector-space bases
$B_1,\ldots,B_{c(p)}$ of $V$, their multiset union is an additive basis.
Writing $B_i=(b_{i1},\ldots,b_{in})$, the assertion is
\[
 \left\{\sum_{i=1}^{c(p)}\sum_{j=1}^n
                 \epsilon_{ij}b_{ij}:\epsilon_{ij}\in\{0,1\}\right\}=V.
\]
The bases need not be distinct and may share vectors. All sums and the
coefficients $0,1$ are interpreted in $\mathbb F_p$.

The constant must be independent of dimension and independent of the
chosen bases. A bound which increases with $n$ would not settle this
question. The source also suggests the stronger candidate $c(p)=p$;
this is distinguished from the principal existence assertion, which
allows any finite value depending on $p$. Saying that the union
``contains'' an additive basis is equivalent to saying that the union
itself has the subset-sum property, since additional occurrences can
always be left unused.
\subsection{Short English statement}
For each prime, do a fixed number of arbitrary vector-space bases supply every vector as a sum using each available occurrence at most once?
\subsection{Sources}
\begin{itemize}
\item[S1] ULAM.AI and the UnsolvedMath Contributors, supplied problems.json, record 3060 (OPG-150), OpenGarden collection. Catalogue identity and source transcription; CC BY 4.0.
\url{https://huggingface.co/datasets/ulamai/UnsolvedMath}.
\item[S2] Open Problem Garden, The additive basis conjecture. Original problem and discussion; the mathematical formulation below is expanded from the supplied source record.
\url{http://www.openproblemgarden.org/op/the_additive_basis_conjecture}.
\end{itemize}
\end{document}
